\section{总结与延伸}

\subsection{讲者的核心总结}

\begin{figure}[H]
\centering
\includegraphics[width=0.95\textwidth]{slides-images/slide-076.jpg}
\caption{Lecture 13 总结：覆盖了生成模型家族树的一半——显式密度（自回归、VAE）和直接隐式密度（GAN）\protect\footnotemark}
\end{figure}
\footnotetext{Slides 第77页。}

Justin Johnson 在课程结尾总结了本节课的三大主题：

\begin{enumerate}
    \item \textbf{生成模型的概率基础}：判别式 vs 生成式 vs 条件生成式，归一化约束的含义
    \item \textbf{显式密度方法}：自回归模型（可精确计算密度），VAE（可计算近似密度下界）
    \item \textbf{隐式密度方法}：GAN（直接采样，但不可计算密度）
\end{enumerate}

下一节课将覆盖家族树的另一半——\textbf{扩散模型}，它属于隐式密度的迭代采样方法。

\subsection{全课知识图谱}

\begin{center}
\begin{tikzpicture}[
    node distance=0.6cm,
    every node/.style={draw, rounded corners, minimum width=2.8cm, minimum height=0.6cm, font=\small},
    level1/.style={fill=blue!10},
    level2/.style={fill=yellow!15},
    level3/.style={fill=green!10}
]
\node[level1] (gen) {生成模型 $p(X)$};
\node[level2, below left=0.8cm and 0.5cm of gen] (explicit) {显式密度};
\node[level2, below right=0.8cm and 0.5cm of gen] (implicit) {隐式密度};
\node[level3, below left=0.6cm and -0.5cm of explicit] (ar) {自回归};
\node[level3, below right=0.6cm and -0.5cm of explicit] (vae) {VAE};
\node[level3, below left=0.6cm and -0.5cm of implicit] (gan) {GAN};
\node[level3, below right=0.6cm and -0.5cm of implicit] (diff) {扩散模型};

\draw[->, thick] (gen) -- (explicit);
\draw[->, thick] (gen) -- (implicit);
\draw[->, thick] (explicit) -- (ar);
\draw[->, thick] (explicit) -- (vae);
\draw[->, thick] (implicit) -- (gan);
\draw[->, thick] (implicit) -- (diff);

\node[draw=none, below=0.1cm of ar, font=\scriptsize, text width=2cm, align=center] {精确密度\\链式分解};
\node[draw=none, below=0.1cm of vae, font=\scriptsize, text width=2cm, align=center] {近似密度\\ELBO};
\node[draw=none, below=0.1cm of gan, font=\scriptsize, text width=2cm, align=center] {直接采样\\对抗训练};
\node[draw=none, below=0.1cm of diff, font=\scriptsize, text width=2cm, align=center] {迭代采样\\去噪过程};
\end{tikzpicture}
\end{center}

\subsection{关键 Takeaways}

\begin{importantbox}{五条核心原则}
\begin{enumerate}
    \item \textbf{归一化约束是生成模型的核心}：不同概率模型的本质区别在于什么东西在竞争概率质量
    \item \textbf{条件生成模型最有实用价值}：$p(X|Y)$ 才是真正有用的模型，$p(X)$ 更多是理论工具
    \item \textbf{MLE 是显式密度方法的统一框架}：自回归和 VAE 都在（直接或间接地）最大化数据的对数似然
    \item \textbf{每种方法都有其权衡}：自回归（精确但慢）、VAE（有 latent space 但模糊）、GAN（清晰但难训练）
    \item \textbf{扩散模型最终统一了质量与稳定性}：解决了 GAN 的训练问题，同时达到甚至超越 GAN 的生成质量
\end{enumerate}
\end{importantbox}

\subsection{拓展阅读}

\begin{itemize}
    \item Kingma \& Welling, ``Auto-Encoding Variational Bayes'' (2013): \url{https://arxiv.org/abs/1312.6114} --- VAE 原始论文
    \item Goodfellow et al., ``Generative Adversarial Nets'' (2014): \url{https://arxiv.org/abs/1406.2661} --- GAN 原始论文
    \item van den Oord et al., ``Pixel Recurrent Neural Networks'' (2016): \url{https://arxiv.org/abs/1601.06759} --- PixelRNN
    \item Karras et al., ``A Style-Based Generator Architecture for GANs'' (2019): \url{https://arxiv.org/abs/1812.04948} --- StyleGAN
    \item Lilian Weng's Blog, ``From Autoencoder to Beta-VAE'': \url{https://lilianweng.github.io/posts/2018-08-12-vae/} --- VAE 详解
\end{itemize}

\end{document}
